\textcolor{red}{\textbf{La primera respuesta generada por la IA produjo un bloque \texttt{flowchart LR} sintácticamente válido, que renderizó sin errores en mermaid.live y respetaba la mayoría de las reglas de la plantilla: una única asociación actor--caso de uso, estereotipos con comillas latinas y \texttt{classDef}/\texttt{class} para el actor. No obstante, tras auditar el código frente a la semántica UML, a los criterios de aceptación de US-7 y al glosario común, se aplicaron las siguientes correcciones y mejoras:
\begin{itemize}[leftmargin=*,nosep]
\item \textbf{Relación «extend» con la dirección invertida:} La IA escribió «Pagar con Bizum» como origen de la flecha «extend» hacia «Autorización en la app del banco». En una relación «extend» la flecha discontinua debe partir del caso de uso opcional y apuntar al caso base. Además, el caso de uso opcional que realmente existe en US-7 es el fallo de autorización, y no la autorización en sí.
\item \textbf{Paso obligatorio modelado como opcional:} La autorización en la app del banco es un paso necesario para completar el pago con Bizum, no un flujo condicional. Se corrigió como caso incluido, «Autorizar pago en la app del banco», y se añadió «Mostrar error de autorización denegada o expirada» como la verdadera extensión «extend», coherente con la rama de pago denegado del diagrama de secuencia de US-7.
\item \textbf{Criterios de aceptación incompletos:} El criterio exige que la pasarela solicite el teléfono del comprador, paso que la IA no modelaba, y tampoco aparecía la confirmación del pedido. Se añadieron los casos incluidos «Solicitar teléfono Bizum» y «Confirmar pedido», y se mantuvo «Redirigir a la pasarela oficial de Bizum» como caso incluido.
\item \textbf{Actor, identificadores y límite inconsistentes con el proyecto:} La IA usó el actor «Clienta Habitual» en lugar de \texttt{Cliente} (glosario común), el identificador \texttt{US7} para un caso de uso (confunde un caso con el código de la historia) y el límite genérico «E-commerce Boutique Moda». Se unificó el actor, se usaron identificadores descriptivos y el límite pasó a titularse «Atelier Jaén - Checkout con Bizum».
\end{itemize}}}

Describid aquí los errores \textbf{reales} de la primera respuesta de la IA para el caso de uso de US-7 y cómo los corrigió el equipo. Puntos que hay que comprobar:
\begin{itemize}[leftmargin=*,nosep]
\item Más de una asociación actor-caso de uso (solo se permite una).
\item Flecha include en sentido contrario (debe ir del caso base al incluido).
\item Flecha extend en sentido contrario (debe ir del opcional al base) o no discontinua.
\item Estereotipos escritos con < > o << >> en lugar de « ».
\item Falta classDef/class del actor o falta el icono fa:fa-user.
\item ID del subgraph repetido como ID de nodo; paréntesis en etiquetas.
\item Casos de uso inventados que no aparecen en los criterios de aceptación, o criterios sin reflejar.
\end{itemize}
